\chapter{Detektoren und Teilchenidentifikation}

\section{Ziel dieses Kapitels}

Dieses Kapitel behandelt, wie Teilchen in Experimenten nachgewiesen und identifiziert
werden.
In der Kern- und Teilchenphysik sieht man Teilchen fast nie direkt.
Man misst stattdessen ihre Wechselwirkungen mit Materie.

\begin{notebox}
Dieses Kapitel behandelt die Stichwoerter:
\begin{itemize}
    \item Teilchendetektoren,
    \item Ionisation,
    \item Energieverlust geladener Teilchen,
    \item Bethe-Bloch-Formel,
    \item Spurdetektoren,
    \item Nebelkammer,
    \item Blasenkammer,
    \item Time Projection Chamber,
    \item Szintillatoren,
    \item Cherenkov-Strahlung,
    \item Kalorimeter,
    \item elektromagnetische Schauer,
    \item hadronische Schauer,
    \item Myondetektoren,
    \item Neutrinodetektion,
    \item Teilchenidentifikation,
    \item Event Display.
\end{itemize}
\end{notebox}

\section{Grundidee eines Teilchendetektors}

Ein Detektor misst nicht das Teilchen selbst, sondern die Spuren seiner Wechselwirkung
mit Materie.

\begin{definitionbox}
Ein \emph{Teilchendetektor} ist ein Messgeraet, das Teilchen durch ihre Wechselwirkung
mit Materie nachweist.
\end{definitionbox}

\begin{conceptbox}
Teilchen koennen in Materie verschiedene Signale erzeugen:
\begin{itemize}
    \item Ionisation,
    \item Anregung von Atomen oder Molekuelen,
    \item Szintillationslicht,
    \item Cherenkov-Licht,
    \item elektromagnetische Schauer,
    \item hadronische Schauer,
    \item Spuren in Gasen, Fluessigkeiten oder Halbleitern.
\end{itemize}
\end{conceptbox}

\begin{examtipbox}
Wichtiger Merksatz:
Detektoren messen Energie, Impuls, Ort, Zeit und Teilchenart nicht direkt auf einmal.
Man kombiniert verschiedene Detektorschichten, um diese Groessen gemeinsam zu rekonstruieren.
\end{examtipbox}

\section{Welche Teilchen sind direkt sichtbar?}

Nicht alle Teilchen hinterlassen auf gleiche Weise Signale im Detektor.

\begin{conceptbox}
Geladene Teilchen sind relativ einfach nachzuweisen, weil sie Materie ionisieren.

Neutrale Teilchen sind schwieriger nachzuweisen.
Sie hinterlassen keine direkte Ionisationsspur.
Man erkennt sie ueber ihre Wechselwirkungen, Zerfaelle oder fehlende Energie und
fehlenden Impuls.
\end{conceptbox}

\begin{notebox}
Typische Detektorsignaturen:
\begin{itemize}
    \item Elektronen: Spur und elektromagnetischer Schauer.
    \item Photonen: kein Track vor dem Kalorimeter, aber elektromagnetischer Schauer.
    \item Myonen: durchdringen viele Detektorschichten und erreichen Myonkammern.
    \item Hadronen: hadronische Schauer.
    \item Neutrinos: meist nur indirekt ueber fehlende Energie oder spezielle Nachweisreaktionen.
\end{itemize}
\end{notebox}

\section{Ionisation}

Geladene Teilchen verlieren in Materie Energie, indem sie Atome ionisieren oder anregen.

\begin{definitionbox}
\emph{Ionisation} bedeutet, dass ein Atom oder Molekuel ein Elektron verliert und
dadurch ein Ion entsteht.
\end{definitionbox}

\begin{conceptbox}
Ein geladenes Teilchen bewegt sich durch Materie und wechselwirkt elektromagnetisch
mit den Elektronen der Atome.
Dabei kann es:
\begin{itemize}
    \item Elektronen herausloesen,
    \item Atome anregen,
    \item Energie an das Medium abgeben.
\end{itemize}
\end{conceptbox}

\begin{definitionbox}
Der \emph{spezifische Energieverlust} ist die pro Wegstrecke abgegebene Energie:
\[
    -\frac{\dd E}{\dd x}.
\]
\end{definitionbox}

\begin{examtipbox}
Das Minuszeichen in
\[
    -\frac{\dd E}{\dd x}
\]
zeigt:
Das Teilchen verliert Energie, waehrend es Materie durchlaeuft.
\end{examtipbox}

\section{Bethe-Bloch-Formel}

Die Bethe-Bloch-Formel beschreibt den mittleren Energieverlust schwerer geladener
Teilchen in Materie.

\begin{definitionbox}
Die \emph{Bethe-Bloch-Formel} beschreibt den mittleren Energieverlust geladener
Teilchen durch Ionisation und Anregung in Materie.
\end{definitionbox}

\begin{formulabox}
Qualitativ beschreibt die Bethe-Bloch-Formel:
\[
    -\frac{\dd E}{\dd x}
    =
    f(z,\beta,\gamma,Z,A,I).
\]
Dabei ist:
\begin{itemize}
    \item $z$ die Ladungszahl des Projektils,
    \item $\beta=v/c$,
    \item $\gamma$ der Lorentzfaktor,
    \item $Z$ die Kernladungszahl des Materials,
    \item $A$ die molare Masse des Materials,
    \item $I$ die mittlere Anregungsenergie des Materials.
\end{itemize}
\end{formulabox}

\begin{conceptbox}
Die wichtigste qualitative Abhaengigkeit ist:
\[
    -\frac{\dd E}{\dd x}
    \propto
    \frac{z^2}{\beta^2}
\]
im nichtrelativistischen Bereich.

Langsame geladene Teilchen ionisieren daher besonders stark.
\end{conceptbox}

\begin{notebox}
Die Vorlesungsfolie zur Bethe-Bloch-Formel verwendet die Materiedicke in
\[
    \mathrm{g/cm^2}
\]
und enthaelt die Parameter $Z$, $A$, $I$, $\beta$, $\gamma$ und die Projektil-Ladung
$z$.
\end{notebox}

\begin{examtipbox}
Pruefungsrelevant ist meist nicht das Auswendiglernen der vollen Bethe-Bloch-Formel,
sondern das Verstaendnis:
\begin{itemize}
    \item geladene Teilchen verlieren Energie durch Ionisation,
    \item langsame Teilchen ionisieren staerker,
    \item der Energieverlust haengt von Teilchenladung und Material ab,
    \item Bethe-Bloch hilft bei Teilchenidentifikation.
\end{itemize}
\end{examtipbox}

\section{Minimum ionizing particles}

Viele relativistische geladene Teilchen haben einen minimalen Energieverlust.
Man nennt sie minimum ionizing particles.

\begin{definitionbox}
Ein \emph{minimum ionizing particle}, kurz MIP, ist ein geladenes Teilchen, dessen
Energieverlust in Materie nahe dem Minimum der Bethe-Bloch-Kurve liegt.
\end{definitionbox}

\begin{conceptbox}
MIPs sind wichtig, weil sie in vielen Detektoren ein typisches Standardsignal liefern.
Myonen in Hochenergieexperimenten verhalten sich oft naeherungsweise wie MIPs.
\end{conceptbox}

\begin{notebox}
Ein Teilchen mit sehr kleinem Energieverlust kann viele Detektorschichten durchdringen.
Das ist ein Grund, warum Myonen oft in aeusseren Myonkammern nachgewiesen werden.
\end{notebox}

\section{Bragg-Peak}

Der Energieverlust geladener Teilchen steigt am Ende ihrer Bahn oft stark an.

\begin{definitionbox}
Der \emph{Bragg-Peak} ist das Maximum der Energieabgabe eines geladenen Teilchens
kurz vor dem Ende seiner Reichweite in Materie.
\end{definitionbox}

\begin{conceptbox}
Wenn ein geladenes Teilchen langsamer wird, steigt sein Energieverlust pro Wegstrecke
zunaechst an.

Kurz vor dem Stillstand gibt es ein Maximum der Ionisation.
Dieses Maximum ist der Bragg-Peak.
\end{conceptbox}

\begin{examtipbox}
Der Bragg-Peak ist besonders wichtig fuer schwere geladene Teilchen wie Alpha-Teilchen
oder Ionen.

Er ist auch wichtig in der Teilchentherapie, weil die Energie lokal am Ende der
Reichweite deponiert werden kann.
\end{examtipbox}

\section{Spurdetektoren}

Spurdetektoren messen die Bahn geladener Teilchen.

\begin{definitionbox}
Ein \emph{Spurdetektor} misst die Positionen, an denen ein geladenes Teilchen den
Detektor durchquert.
Aus mehreren Messpunkten rekonstruiert man die Teilchenspur.
\end{definitionbox}

\begin{conceptbox}
In einem Magnetfeld sind die Spuren geladener Teilchen gekruemmt.
Aus der Kruemmung kann man den Impuls bestimmen.

Je groesser der Impuls, desto geringer ist die Kruemmung.
\end{conceptbox}

\begin{formulabox}
Fuer ein Teilchen mit Ladung $q$ in einem Magnetfeld gilt:
\[
    p = qBr.
\]
In praktischen Einheiten:
\[
    p[\mathrm{GeV}/c]
    \approx
    0.3\,B[\mathrm{T}]\,r[\mathrm{m}].
\]
\end{formulabox}

\begin{examtipbox}
Spurdetektoren liefern:
\begin{itemize}
    \item Ort der Teilchenspur,
    \item Ladungsvorzeichen aus der Kruemmungsrichtung,
    \item Impuls aus dem Kruemmungsradius,
    \item Vertices aus Schnittpunkten mehrerer Spuren.
\end{itemize}
\end{examtipbox}

\section{Nebelkammer und Blasenkammer}

Fruehe Detektoren machten Teilchenspuren direkt sichtbar.

\begin{definitionbox}
Eine \emph{Nebelkammer} macht Teilchenspuren durch Kondensation entlang ionisierter
Bahnen sichtbar.
\end{definitionbox}

\begin{definitionbox}
Eine \emph{Blasenkammer} macht Teilchenspuren durch Blasenbildung in einer
ueberhitzten Fluessigkeit sichtbar.
\end{definitionbox}

\begin{conceptbox}
Geladene Teilchen ionisieren das Medium.
Entlang dieser Ionisationsspur bilden sich sichtbare Tropfen oder Blasen.

Dadurch entstehen fotografierbare Spuren.
\end{conceptbox}

\begin{notebox}
Die historischen Detektorabbildungen in den Vorlesungsfolien zeigen Teilchenspuren,
z. B. in Zusammenhang mit Positron-Entdeckung und Blasenkammer-/Spuraufnahmen.
\end{notebox}

\begin{examtipbox}
Nebelkammer und Blasenkammer sind heute vor allem historisch und konzeptionell wichtig.
Sie zeigen sehr anschaulich:
Geladene Teilchen hinterlassen Spuren durch Ionisation.
\end{examtipbox}

\section{Time Projection Chamber}

Eine Time Projection Chamber, kurz TPC, ist ein moderner Spurdetektor.

\begin{definitionbox}
Eine \emph{Time Projection Chamber} ist ein gasgefuellter Spurdetektor, der die
dreidimensionale Spur geladener Teilchen rekonstruiert.
\end{definitionbox}

\begin{conceptbox}
Ein geladenes Teilchen ionisiert das Gas.
Die freigesetzten Elektronen driften in einem elektrischen Feld zu einer Ausleseebene.

Aus der Position an der Ausleseebene und der Driftzeit rekonstruiert man die
dreidimensionale Spur.
\end{conceptbox}

\begin{notebox}
Die Vorlesungsabbildungen enthalten eine TPC-Abbildung.
Eine TPC verbindet Ionisationsmessung mit raeumlicher Spurrekonstruktion.
\end{notebox}

\begin{examtipbox}
TPC:
\begin{itemize}
    \item misst Spuren geladener Teilchen,
    \item nutzt Ionisation im Gas,
    \item Driftzeit liefert die dritte Koordinate,
    \item kann auch Energieverlustinformation fuer Teilchenidentifikation liefern.
\end{itemize}
\end{examtipbox}

\section{Halbleiterdetektoren}

Moderne Experimente verwenden haeufig Halbleiterdetektoren nahe am Kollisionspunkt.

\begin{definitionbox}
Ein \emph{Halbleiterdetektor} weist Teilchen nach, indem geladene Teilchen in einem
Halbleiter Elektron-Loch-Paare erzeugen.
\end{definitionbox}

\begin{conceptbox}
Halbleiterdetektoren besitzen sehr gute Ortsaufloesung.
Sie sind deshalb ideal fuer innere Spurdetektoren.

Man verwendet sie, um:
\begin{itemize}
    \item Primaervertices zu rekonstruieren,
    \item Sekundaervertices zu finden,
    \item kurzlebige Teilchen indirekt nachzuweisen,
    \item Impulse geladener Teilchen zu messen.
\end{itemize}
\end{conceptbox}

\begin{examtipbox}
Sekundaervertices sind wichtig fuer Teilchen mit messbarer Lebensdauer,
zum Beispiel Hadronen mit schweren Quarks.
\end{examtipbox}

\section{Szintillatoren}

Szintillatoren wandeln Teilchenenergie in Licht um.

\begin{definitionbox}
Ein \emph{Szintillator} ist ein Material, das bei Durchgang ionisierender Strahlung
Licht aussendet.
\end{definitionbox}

\begin{conceptbox}
Ein Teilchen regt Atome oder Molekuele im Szintillatormaterial an.
Bei der Abregung wird Licht ausgesendet.
Dieses Licht kann mit Photomultipliern oder Photodetektoren gemessen werden.
\end{conceptbox}

\begin{notebox}
Szintillatoren werden haeufig fuer Trigger, Zeitmessung und Kalorimetrie verwendet.
\end{notebox}

\begin{examtipbox}
Szintillatoren sind schnelle Detektoren.
Sie sind daher besonders nuetzlich fuer Zeitmessung und Triggerentscheidungen.
\end{examtipbox}

\section{Cherenkov-Strahlung}

Cherenkov-Strahlung entsteht, wenn ein geladenes Teilchen in einem Medium schneller
ist als Licht in diesem Medium.

\begin{definitionbox}
\emph{Cherenkov-Strahlung} ist Licht, das ausgesendet wird, wenn ein geladenes
Teilchen ein Medium mit einer Geschwindigkeit durchquert, die groesser ist als die
Lichtgeschwindigkeit in diesem Medium.
\end{definitionbox}

\begin{formulabox}
Die Bedingung fuer Cherenkov-Strahlung ist:
\[
    v > \frac{c}{n},
\]
oder
\[
    \beta n > 1.
\]
Dabei ist:
\begin{itemize}
    \item $v$ die Teilchengeschwindigkeit,
    \item $c$ die Lichtgeschwindigkeit im Vakuum,
    \item $n$ der Brechungsindex des Mediums,
    \item $\beta=v/c$.
\end{itemize}
\end{formulabox}

\begin{formulabox}
Der Cherenkov-Winkel erfuellt:
\[
    \cos\theta_C = \frac{1}{\beta n}.
\]
\end{formulabox}

\begin{conceptbox}
Aus dem Cherenkov-Winkel kann man die Geschwindigkeit des Teilchens bestimmen.
Kombiniert man die Geschwindigkeit mit dem Impuls aus dem Spurdetektor, kann man
die Masse und damit die Teilchenart bestimmen.
\end{conceptbox}

\begin{examtipbox}
Cherenkov-Detektoren sind Teilchenidentifikationsdetektoren.
Sie liefern Geschwindigkeitsinformation.
Zusammen mit dem Impuls kann man Teilchen unterscheiden.
\end{examtipbox}

\section{Kalorimeter}

Kalorimeter messen die Energie von Teilchen, indem sie diese in Materie stoppen.

\begin{definitionbox}
Ein \emph{Kalorimeter} ist ein Detektor, der die Energie eines Teilchens misst,
indem das Teilchen im Detektormaterial einen Schauer erzeugt und seine Energie
deponiert.
\end{definitionbox}

\begin{conceptbox}
Es gibt zwei wichtige Typen:
\begin{itemize}
    \item elektromagnetische Kalorimeter,
    \item hadronische Kalorimeter.
\end{itemize}
\end{conceptbox}

\begin{definitionbox}
Ein \emph{elektromagnetisches Kalorimeter} misst vor allem Elektronen, Positronen
und Photonen.

Ein \emph{hadronisches Kalorimeter} misst Hadronen wie Protonen, Neutronen, Pionen
und Kaonen.
\end{definitionbox}

\begin{examtipbox}
Kalorimeter liefern Energieinformation.
Spurdetektoren liefern Impuls- und Ortsinformation.
Beides zusammen ist fuer Teilchenidentifikation sehr wichtig.
\end{examtipbox}

\section{Elektromagnetische Schauer}

Elektronen, Positronen und Photonen erzeugen elektromagnetische Schauer.

\begin{conceptbox}
Ein hochenergetisches Elektron verliert Energie durch Bremsstrahlung:
\[
    e^- \rightarrow e^- + \gamma.
\]

Ein hochenergetisches Photon kann Paarbildung machen:
\[
    \gamma \rightarrow e^- + e^+.
\]

Diese Prozesse wiederholen sich und erzeugen einen elektromagnetischen Schauer.
\end{conceptbox}

\begin{definitionbox}
Ein \emph{elektromagnetischer Schauer} ist eine Kaskade aus Elektronen, Positronen
und Photonen, die durch Bremsstrahlung und Paarbildung entsteht.
\end{definitionbox}

\begin{examtipbox}
Elektronen und Photonen werden im elektromagnetischen Kalorimeter gestoppt.

Unterschied:
\begin{itemize}
    \item Elektron: hat eine Spur im inneren Detektor und einen elektromagnetischen Schauer.
    \item Photon: hat normalerweise keine Spur vor dem Kalorimeter, aber einen elektromagnetischen Schauer.
\end{itemize}
\end{examtipbox}

\section{Hadronische Schauer}

Hadronen wechselwirken stark mit Atomkernen im Detektormaterial.

\begin{definitionbox}
Ein \emph{hadronischer Schauer} ist eine Kaskade aus Hadronen und Sekundaerteilchen,
die durch starke Wechselwirkungen in Materie entsteht.
\end{definitionbox}

\begin{conceptbox}
Hadronische Schauer sind im Allgemeinen komplexer als elektromagnetische Schauer,
weil starke Wechselwirkungen, Kernaufbrueche und viele Sekundaerteilchen beteiligt sind.
\end{conceptbox}

\begin{examtipbox}
Hadronen werden vor allem im hadronischen Kalorimeter nachgewiesen.

Typische Hadronen sind:
\begin{itemize}
    \item Protonen,
    \item Neutronen,
    \item geladene und neutrale Pionen,
    \item Kaonen.
\end{itemize}
\end{examtipbox}

\section{Myondetektoren}

Myonen sind relativ durchdringende Teilchen.
Sie erreichen oft die aeusseren Detektorschichten.

\begin{definitionbox}
Ein \emph{Myondetektor} ist ein aeusserer Detektorbereich, der speziell fuer den
Nachweis von Myonen ausgelegt ist.
\end{definitionbox}

\begin{conceptbox}
Myonen sind schwerer als Elektronen und verlieren in Materie oft nur wenig Energie
durch Ionisation.
Sie erzeugen normalerweise keine vollstaendig absorbierten elektromagnetischen oder
hadronischen Schauer im inneren Detektor.

Deshalb koennen sie viele Detektorschichten durchdringen und in den aeusseren
Myonkammern gemessen werden.
\end{conceptbox}

\begin{notebox}
Die LHC-Folien zeigen Myonspuren in einem Event Display, zum Beispiel bei
\[
    pp \rightarrow \mu^+ \mu^- + \ldots
\]
Die Myonen erscheinen als durchgehende Spuren bis in die aeusseren Detektorbereiche.
\end{notebox}

\begin{examtipbox}
Myon-Signatur:
\begin{itemize}
    \item Spur im inneren Detektor,
    \item wenig Energie im Kalorimeter,
    \item Treffer in Myonkammern.
\end{itemize}
\end{examtipbox}

\section{Mehrschichtdetektoren am Collider}

Grosse Collider-Experimente wie ATLAS und CMS bestehen aus mehreren Detektorschichten.

\begin{conceptbox}
Typischer Aufbau von innen nach aussen:
\begin{itemize}
    \item innerer Spurdetektor,
    \item elektromagnetisches Kalorimeter,
    \item hadronisches Kalorimeter,
    \item Myonsystem.
\end{itemize}
\end{conceptbox}

\begin{notebox}
Die LHC-Folien zeigen ATLAS und CMS als grosse Mehrzweckdetektoren.
ATLAS ist sehr gross und besitzt eine zylindrische, mehrschichtige Detektorgeometrie.
\end{notebox}

\begin{conceptbox}
Die Schichten sind so angeordnet, dass verschiedene Teilchenarten unterschiedliche
Signaturen hinterlassen.

Dadurch kann man aus dem Muster der Detektorsignale auf die Teilchenart schliessen.
\end{conceptbox}

\section{Teilchenidentifikation im Mehrschichtdetektor}

\begin{notebox}
Typische Signaturen:
\begin{itemize}
    \item Photon: elektromagnetischer Kalorimeterschauer ohne zugehoerige geladene Spur.
    \item Elektron: geladene Spur plus elektromagnetischer Kalorimeterschauer.
    \item Myon: geladene Spur plus Treffer im Myonsystem.
    \item Geladenes Hadron: geladene Spur plus hadronischer Kalorimeterschauer.
    \item Neutrales Hadron: hadronischer Kalorimeterschauer ohne geladene Spur.
    \item Neutrino: fehlende transversale Energie beziehungsweise fehlender Impuls.
\end{itemize}
\end{notebox}

\begin{examtipbox}
Die Teilchenidentifikation erfolgt durch Kombination:
\[
    \text{Spur}
    +
    \text{Kalorimeterenergie}
    +
    \text{Myonsystem}
    +
    \text{fehlende Energie}.
\]
\end{examtipbox}

\section{Neutrinodetektion}

Neutrinos wechselwirken nur schwach.
Deshalb sind sie extrem schwer nachzuweisen.

\begin{definitionbox}
Ein \emph{Neutrinodetektor} ist ein sehr grosser Detektor, der seltene schwache
Wechselwirkungen von Neutrinos mit Materie nachweist.
\end{definitionbox}

\begin{conceptbox}
Neutrinos sind elektrisch neutral und wechselwirken nicht elektromagnetisch.
Sie ionisieren Materie nicht direkt.

Man weist sie nach, indem sie gelegentlich eine schwache Wechselwirkung mit einem
Atomkern oder Elektron ausloesen.
Dabei entstehen geladene Teilchen, die dann detektiert werden koennen.
\end{conceptbox}

\begin{notebox}
Die Vorlesungsabbildungen enthalten eine Super-Kamiokande-Abbildung.
Super-Kamiokande ist ein grosser Wasser-Cherenkov-Detektor.
Geladene Teilchen, die durch Neutrinowechselwirkungen entstehen, koennen Cherenkov-Licht
im Wasser erzeugen.
\end{notebox}

\begin{examtipbox}
Neutrinos sieht man nicht direkt.
Man sieht geladene Sekundaerteilchen aus Neutrinowechselwirkungen oder rekonstruiert
fehlende Energie und fehlenden Impuls.
\end{examtipbox}

\section{Fehlende Energie und fehlender Impuls}

In Collider-Experimenten werden Neutrinos haeufig indirekt rekonstruiert.

\begin{definitionbox}
\emph{Fehlende transversale Energie} ist ein Hinweis darauf, dass unsichtbare Teilchen
wie Neutrinos den Detektor verlassen haben.
\end{definitionbox}

\begin{conceptbox}
In einem Collider ist der Anfangsimpuls quer zur Strahlrichtung naeherungsweise null.

Wenn die sichtbaren Endprodukte insgesamt einen von null verschiedenen transversalen
Impuls haben, muss ein unsichtbares Teilchen den fehlenden transversalen Impuls
getragen haben.
\end{conceptbox}

\begin{formulabox}
Die fehlende transversale Impulsbilanz ist grob:
\[
    \vec p_{\mathrm{T}}^{\,\mathrm{miss}}
    =
    -
    \sum_{\mathrm{sichtbar}}
    \vec p_{\mathrm{T},i}.
\]
\end{formulabox}

\begin{examtipbox}
Fehlende transversale Energie ist typisch fuer Ereignisse mit Neutrinos oder anderen
unsichtbaren Teilchen.
\end{examtipbox}

\section{Trigger}

Collider erzeugen extrem viele Ereignisse.
Nicht alle koennen gespeichert werden.

\begin{definitionbox}
Ein \emph{Trigger} ist ein System, das schnell entscheidet, welche Ereignisse
gespeichert und genauer analysiert werden.
\end{definitionbox}

\begin{conceptbox}
Ein Trigger sucht nach interessanten Signaturen, zum Beispiel:
\begin{itemize}
    \item hochenergetische Elektronen,
    \item Myonen,
    \item Photonen,
    \item Jets,
    \item grosse fehlende transversale Energie.
\end{itemize}
\end{conceptbox}

\begin{notebox}
Ohne Trigger waeren moderne Collider-Experimente nicht auswertbar, weil die Rohdatenrate
viel zu gross ist.
\end{notebox}

\section{Event Display}

Ein Event Display ist eine grafische Darstellung eines einzelnen Kollisionsereignisses.

\begin{definitionbox}
Ein \emph{Event Display} ist eine Visualisierung der in einem Detektor gemessenen
Spuren, Energiedepositionen und Treffer eines einzelnen Ereignisses.
\end{definitionbox}

\begin{conceptbox}
In einem Event Display erkennt man zum Beispiel:
\begin{itemize}
    \item gekruemmte Spuren geladener Teilchen,
    \item Energiecluster in Kalorimetern,
    \item Myonspuren in aeusseren Detektorschichten,
    \item Jets als gebuendelte Teilchenbuendel,
    \item fehlende Energie aus Impulsungleichgewicht.
\end{itemize}
\end{conceptbox}

\begin{notebox}
Die LHC-Folien zeigen Event Displays, unter anderem mit Myonen.
Solche Darstellungen sind wichtig, um die Rekonstruktion von Ereignissen anschaulich
zu verstehen.
\end{notebox}

\section{Zusammenspiel von Messgroessen}

Teilchenidentifikation basiert auf mehreren gemessenen Groessen.

\begin{conceptbox}
Wichtige Messgroessen sind:
\begin{itemize}
    \item Impuls aus Spurkrümmung,
    \item Ladungsvorzeichen aus Krümmungsrichtung,
    \item Energie aus Kalorimetern,
    \item Geschwindigkeit aus Cherenkov- oder Flugzeitmessung,
    \item Energieverlust $\dd E/\dd x$,
    \item Treffer im Myonsystem,
    \item fehlende transversale Energie.
\end{itemize}
\end{conceptbox}

\begin{formulabox}
Wenn Impuls $p$ und Geschwindigkeit $v$ bekannt sind, kann man die Masse bestimmen.
Relativistisch gilt:
\[
    p=\gamma mv.
\]
Mit
\[
    \beta=\frac{v}{c}
\]
und
\[
    \gamma=\frac{1}{\sqrt{1-\beta^2}}
\]
folgt:
\[
    m = \frac{p}{\gamma v}.
\]
\end{formulabox}

\begin{examtipbox}
Teilchenidentifikation bedeutet meistens:
Man misst mehrere Groessen und prueft, welche Teilchenhypothese dazu passt.
\end{examtipbox}

\section{Mini-Zusammenfassung}

\begin{notebox}
Die wichtigsten Punkte dieses Kapitels sind:
\begin{itemize}
    \item Detektoren messen Wechselwirkungen von Teilchen mit Materie.
    \item Geladene Teilchen ionisieren Materie und hinterlassen Spuren.
    \item Der Energieverlust wird durch die Bethe-Bloch-Formel beschrieben.
    \item Spurdetektoren messen Ort, Impuls und Ladungsvorzeichen geladener Teilchen.
    \item Nebelkammern und Blasenkammern machen Ionisationsspuren sichtbar.
    \item Eine TPC rekonstruiert dreidimensionale Spuren aus Driftzeit und Ausleseposition.
    \item Halbleiterdetektoren liefern sehr gute Ortsaufloesung.
    \item Szintillatoren wandeln Teilchenenergie in Licht um.
    \item Cherenkov-Detektoren liefern Geschwindigkeitsinformation.
    \item Kalorimeter messen Energien durch Schauerbildung.
    \item Elektromagnetische Kalorimeter messen Elektronen und Photonen.
    \item Hadronische Kalorimeter messen Hadronen.
    \item Myonen erreichen oft aeussere Myonkammern.
    \item Neutrinos werden indirekt ueber schwache Wechselwirkungen oder fehlende Energie nachgewiesen.
    \item Teilchenidentifikation entsteht aus der Kombination mehrerer Detektorsignale.
\end{itemize}
\end{notebox}

\section{Pruefungscheck}

\begin{examtipbox}
Nach diesem Kapitel solltest du folgende Fragen beantworten koennen:
\begin{enumerate}
    \item Warum sieht man Teilchen nicht direkt?
    \item Welche Signale koennen Teilchen in Materie erzeugen?
    \item Warum sind geladene Teilchen einfacher nachzuweisen als neutrale?
    \item Was beschreibt die Bethe-Bloch-Formel?
    \item Was bedeutet $-\dd E/\dd x$?
    \item Was ist ein minimum ionizing particle?
    \item Was ist der Bragg-Peak?
    \item Wie funktioniert ein Spurdetektor?
    \item Wie erhaelt man aus Spurkrümmung den Impuls?
    \item Was ist der Unterschied zwischen Nebelkammer und Blasenkammer?
    \item Wie funktioniert eine Time Projection Chamber?
    \item Warum sind Halbleiterdetektoren nahe am Kollisionspunkt wichtig?
    \item Was ist ein Szintillator?
    \item Wann entsteht Cherenkov-Strahlung?
    \item Wie lautet die Bedingung fuer Cherenkov-Strahlung?
    \item Was misst ein Kalorimeter?
    \item Was ist ein elektromagnetischer Schauer?
    \item Was ist ein hadronischer Schauer?
    \item Warum erreichen Myonen oft die aeusseren Detektorschichten?
    \item Wie werden Neutrinos nachgewiesen?
    \item Was bedeutet fehlende transversale Energie?
    \item Was ist ein Trigger?
    \item Was zeigt ein Event Display?
    \item Wie kombiniert man Detektorsignale zur Teilchenidentifikation?
\end{enumerate}
\end{examtipbox}